\documentclass[11pt]{article}

\usepackage[margin=1in]{geometry}
\usepackage{enumitem}
\usepackage{hyperref}

\title{Requirements for Creating an iOS IPA Build}
\author{AI Toy Companion Mobile App}
\date{\today}

\begin{document}

\maketitle

\section*{Purpose}

This document summarizes what is required to create an iOS IPA file for the AI Toy Companion mobile application. The current development environment does not include macOS/Xcode access, so the client will need to provide the Apple build environment and Apple developer credentials.

\section*{Required Client Resources}

\begin{itemize}[leftmargin=*]
  \item A Mac computer or Mac build machine with a recent supported version of macOS.
  \item Xcode installed from the Mac App Store or Apple Developer portal.
  \item Apple Developer Program membership for the client organization.
  \item Access to App Store Connect, if the build will be uploaded for TestFlight or App Store review.
  \item iOS bundle identifier, for example: \texttt{com.client.aitoycompanion}.
  \item Signing certificate and provisioning profile, or permission for Xcode to manage signing automatically.
  \item Target iPhone/iPad test devices registered in the Apple Developer account if creating an ad hoc build.
\end{itemize}

\section*{Project Build Requirements}

\begin{itemize}[leftmargin=*]
  \item Node.js and npm installed on the Mac.
  \item Ruby/Bundler and CocoaPods available for installing iOS native dependencies.
  \item Project dependencies installed using \texttt{npm install}.
  \item iOS pods installed using:
  \begin{verbatim}
cd ios
bundle install
bundle exec pod install
  \end{verbatim}
  \item Required app configuration values, API endpoints, and environment values provided before archiving.
\end{itemize}

\section*{IPA Creation Process}

\begin{enumerate}[leftmargin=*]
  \item Open the iOS workspace in Xcode: \texttt{ios/*.xcworkspace}.
  \item Select the correct app target and development team.
  \item Confirm bundle identifier, version number, build number, permissions, and signing settings.
  \item Build and test the app on a real iPhone.
  \item Create an archive using Xcode: \texttt{Product > Archive}.
  \item Export the archive as an IPA for the required distribution type:
  \begin{itemize}
    \item Development build
    \item Ad hoc build
    \item TestFlight/App Store build
    \item Enterprise build, if the client has enterprise distribution
  \end{itemize}
\end{enumerate}

\section*{Notes}

The IPA cannot be generated from a Windows-only environment because Apple requires Xcode and Apple code signing tools for iOS builds. Once the client provides Mac/Xcode access and signing credentials, the development team can prepare, archive, and export the IPA.

\end{document}
